\section{Discussion: interfaces and the division of proof work}
\label{sec:discussion}

\subsection{What the mathematical comparison reveals}

The downstream operation determines what a preservation statement must
expose. Transferring an accuracy guarantee requires the relevant event
probability; equality of output laws supports all such transfers.
A named-coordinate equation uses the coordinate map. Ambient
differentiation uses local information about a field. A sharp norm
estimate uses the normalization and weights of the preceding bounds.
Exact data, laws, estimates, and local differential information can
therefore appear together in one development.

The most informative transitions occur when the next operation changes
these requirements. The common-base proof preserves a distribution
through reindexing, then execution on each supplied bank. Mahler
identifies form-valued fields for differentiation, then uses tangential
agreement for boundary integration. Trace-cone reconstruction obtains
ideal order from addition, then uses convergence and scaling to identify
the ideals of a particular trace. Following the next operation explains
why each interface has its particular fields.

Implementation details sharpen the comparison. Encoded width and
numerical attributes answer different complexity questions in the
compiler. The graph restriction identity restores the counting norm
when local assignment estimates are summed. Image filters extend a
carrier-indexed convergence property to arbitrary nets. These choices
expose mathematical data through lemmas that later constructions can
reuse. The compiled ideal-composition example in
Section~\ref{sec:upstream-reuse} demonstrates this reuse directly.

Two cases show how a local calculation controls a global estimate.
The graph's total-two table isolates an exceptional direction whose
aggregate energy is paid for by switch gains. In Crouzeix, rectangular
compressions and two orders of polynomial multiplication produce
compatible energy bounds. The final scalar cancellation is short because
the earlier constructions have supplied its exact constants and signs.

These patterns have substantial antecedents. Verified compilers establish
semantic preservation between source and implementation
\citep{Leroy2009,CakeML2014}; probabilistic program logics reason about
distributions and observations \citep{CertiCrypt2009,EasyCrypt2011}.
Algebraic hierarchy design and specification-driven collaboration make
interfaces part of mathematical construction
\citep{GarillotEtAl2009,Baanen2022,CommelinTopaz2024}. The present
comparison identifies what those interfaces preserve in the selected
proofs and how their composition supplies a later calculation's inputs.

\subsection{Tactic use and the division of proof work}
\label{sec:tactics-synthesis}

The corpus survey adds a measured comparison of written proof practice.
After joint subject and proof-length standardization, arithmetic and
algebra tactics remain more prevalent by 22.8 percentage points. Tactic
mode and local claims retain differences of 10.4 and 5.8 points,
respectively. The larger subject-only contrasts in
Figure~\ref{fig:method-shift} also reflect proof-length composition;
Table~\ref{tab:joint-profile} makes that distinction explicit. Arithmetic
emphasis is the clearest persistent difference among these features.

The common-base fold-cost proof shows how such local automation fits
into a semantic argument. Execution identities justify applying the
induction hypothesis to the remaining fold. The safety invariant and
width bound supply its premises. The following consecutive source lines
then apply that hypothesis and close the cost inequality:\footnote{Source:
\href{https://github.com/openai/math/blob/fd4aeeb2ee4fc729c18d98444fed42fd0529eeeb/lean/OAI/Combinatorics/MatroidCounting/CommonBases.lean\#L16880-L16907}{\src{CommonBases.lean}, lines 16880--16907}; excerpt, lines 16905--16907.}
\begin{lstlisting}[caption={Lean source excerpt: fold induction followed by arithmetic discharge.},label={lst:fold-cost}]
have hh := ih z' y' e (hs.2 z' hrun) hw' ht
simp only [List.length_cons,Nat.add_mul,Nat.one_mul]
omega
\end{lstlisting}
Here \src{hs.2} supplies safety after the accumulator update, and
\src{hw'} bounds the remaining list's width. Expanding the successor
length exposes shared product terms. Although the displayed cost bound
contains multiplication, the remaining comparison fits linear
natural-number arithmetic. The semantic and size arguments prepare the
obligation that \src{omega} checks.

\Needspace{14\baselineskip}
The Crouzeix proof supplies a parallel example in analysis. Fourier and
positive-block estimates have already produced bounds with the same
energy $a$ and norm $\gamma$. Their closing scalar argument is:\footnote{Source:
\href{https://github.com/openai/math/blob/fd4aeeb2ee4fc729c18d98444fed42fd0529eeeb/lean/OAI/Analysis/DirectCrouzeix/DomainCore.lean\#L105-L117}{\src{DirectCrouzeix/DomainCore.lean}, lines 105--117}; excerpt, lines 111--117.}
\begin{lstlisting}[caption={Lean source excerpt: the closing Crouzeix energy calculation.},label={lst:crouzeix-close}]
have hb : 2*a ≤ 8*a/γ^2 := by
  simp only [mul_div_assoc] at hPE hQE ⊢
  linarith
have hmul := (le_div_iff₀ (sq_pos_of_pos hγ)).mp hb
change γ ≤ 2
have hg2 : γ^2 ≤ 4 := by nlinarith [hapos]
nlinarith
\end{lstlisting}
The first tactic combines the available energy inequalities. Positivity
of $\gamma$ then permits denominator clearing, and positivity of $a$
permits the final cancellation. The choice of compressions, Fourier
weights, and zero-mode identities creates the relations that these
arithmetic steps certify.

Exact numerical certificates exhibit a further division of work. Their
finite checks verify retained interval data; assembly and soundness
lemmas explain why those checks imply the required estimate. Across
these examples, representation design and local automation meet at an
explicit obligation with proved hypotheses. The survey establishes a
broader method profile, while the case tracing explains mechanisms
within completed source proofs. Testing effects on completion rates,
development effort, or maintenance would require matched tasks and
corresponding outcome measurements.

\subsection{Proof generation, simplification, and mathematical explanation}
\label{sec:generation-explanation}

The comparison identifies differences between completed source corpora.
Mathlib is a collaborative library with varied authoring histories, rather
than a controlled sample of unaided human proofs. These data therefore
cannot establish a distinct AI reasoning process, an advantage in
completion time, or a superior method of doing mathematics. They suggest
a workflow hypothesis: automation may be especially effective when
proved interfaces expose many local obligations in forms handled by
existing tactics. The relevant conditions---available representations,
library coverage, and checkable local goals---can vary within a subject
as well as between subjects. Interface construction remains part of the
work whose cost must be measured.

Obtaining a checked proof and obtaining an economical explanation are
separate objectives. ProofOptimizer reports substantial token
reductions in generated Lean proofs, using Lean to check replacements
against the original theorem statements \citep{GuEtAl2025ProofOptimizer}. Lean Refactor studies proof
length, compilation cost, and compatibility across library versions as
distinct optimization objectives \citep{LuEtAl2026LeanRefactor}. These
results indicate scope for refactoring on the evaluated benchmarks;
they do not measure redundancy in the six
developments examined here. Shorter tactic scripts can also invoke
heavier automation. Source length, proof-checking cost, dependency
structure, and explanatory clarity should consequently be assessed
separately.

Together these observations motivate a research direction: connect
proof generation to compression, reusable interfaces, and mathematical
explanation. A useful interface may amortize its construction cost over
many later results, while a concise final proof may conceal substantial
setup. Claims of improved mathematical practice require evidence about
total effort, generality, downstream reuse, and reader understanding,
alongside correctness. The present findings locate possible mechanisms
and evaluation targets without establishing those improvements.

\subsection{Using the comparison in a new formalization}

Begin with the quantity the next theorem consumes: an event probability,
a coordinate equation, a derivative, or an inequality. Identify what a
proposed representation changes, such as multiplicities, coordinate
order, multiplication order, or the neighborhood of definition. The
preservation statement should connect those changed data to that
specific quantity.

The worked examples make the choice concrete. An event indicator is
sufficient to transfer one accuracy probability, while a full output-law
identity supports reuse for other events. Neighborhood agreement
identifies ambient derivatives; tangential agreement can suffice for the
later integral. A named seam equation needs its coordinate maps.
Proposition~\ref{prop:success-budget} combines several interfaces in one
count: implementation loss, tape uniformity, and deterministic execution
jointly yield the success guarantee.

Next, follow how the hypotheses are supplied for the actual inputs.
Mahler's strip maps construct the analytic data, and Crouzeix's outer
domains provide strict containment and the required estimates. In the
marked-surface construction, one splitting witness supplies both the
diagram and its coordinate identifications. The reusable design consists
of a precise local model, a preservation relation, and a construction of
the data required to apply it.

Finally, inspect the local goals exposed by that design. A width
invariant can reduce a compiler cost step to natural-number arithmetic.
Positive denominators and nonzero factors can prepare an analytic
cancellation. Exact interval data can make a numerical obligation
decidable. The interface proofs assign established automation a concrete
part of the larger mathematical argument to check.

\subsection{Scope and further work}

The six cases were selected for interfaces that could be traced in depth.
Their mathematical comparisons are complemented by the separate uniform
module sample. Joint domain/length controls and directory-family
sensitivity checks sharpen the survey's descriptive interpretation.
A comparison of analogous declaration roles and mathematical tasks would
further separate library composition from choices about proof organization.
Elaborated tactic traces and a held-out extraction-validation cohort
would connect the lexical measurements more directly to executed proof steps.
A prospective comparison could cross AI assistance with availability of
proved representation interfaces on matched tasks, retaining unsuccessful
attempts. Relevant outcomes include time to a checked proof, human
intervention, inference and checking costs, later reuse, and blinded
expert assessments of clarity. This would test the workflow hypothesis
without treating completed-proof length as a causal control.

Source reading identifies the inspected declarations, their hypotheses,
and the constructions connecting them. Compilation of the worked
examples and the three-module trace-cone reuse experiment supplies
additional checking evidence at those targets. Rebuilding the remaining
case endpoints and examining their mathematical arguments would extend
this assessment. Endpoint statements and support results retain their
distinct roles throughout the comparison.

For readers of substantial Lean developments, the practical route is to
locate the next operation, identify what it consumes, and follow the
theorem preserving that information. For formalizers, the cases show how
interface design prepares local proof obligations. The corpus comparison
adds evidence about which written devices are prevalent in this
collection, with arithmetic emphasis persisting under coarse joint
controls. The review brings these two levels together: the mathematical
content of preservation statements and the local proof work they support.

\Needspace{24\baselineskip}
\section*{Data and code availability}

The accompanying reproduction repository (release~1.6.0) provides the
manuscript source, data, and analysis code \citep{GounderArchitectureData2026}.
The source package contains the LaTeX manuscript, reference records,
declaration ledgers, original Lean examples, and Python checks. Five
elementary Lean lemmas import only \src{Lean}. The upstream reuse
experiment imports \src{OAI.Analysis.TraceCone.IdealTransport} and compiles
its three-module dependency chain from unchanged source bytes. Both use
Lean 4.34.1; the latter uses Mathlib revision
\src{d13f23b723b8a846827a245b89c10fc7d3f11612}. Compiler logs retain the
reported axioms. Exact Python checks enumerate tape fibers and Boolean
matrix counts. The README gives reproduction commands and identifies the
external dependencies.

The method-survey directory retains the complete path population, all
800 sampled source modules and their hashes, extraction tables, feature
definitions, bootstrap summaries, and forty inspected validation cases.
The joint comparison additionally retains cell weights, conditional
rates, intervals, and the generated manuscript table. Baseline profiling
requires the pinned Mathlib source checkout; the lexical survey reads
source without elaborating the developments. Public manuscript source
links resolve to the reviewed repository revision.

The source inventory is the complete Git tree of \src{openai/math},
inspected on 8 October 2026.\footnote{Reviewed revision:
\src{fd4aeeb2ee4fc729c18d98444fed42fd0529eeeb}. Public source links refer
to that revision.} It contains 122,875 files with the \src{.lean}
extension, including supporting modules. The package retains 546 source
artifacts, including the \src{OAI} import closures for the six selected
cases; twenty-one detailed records document the additional proof-body
analysis. The separate formalization catalogue lists 173 source records
and 200 main-result bindings \citep{OAICatalog}. These inventories count
different objects and should be read alongside their stated units.

The collection also supplies Comparator challenge configurations
\citep{OAIComparator}. The inspected checker implements statement
comparison, dependency and axiom inspection, kernel replay, and an
optional external checker \citep{ComparatorSource2026}.\footnote{Checker
source revision: \src{ca04cfc72b550331658ec314bf47685281bfd4bf}. This
records the infrastructure source inspected; the collection's
installation instructions do not specify that revision.}
The distinction between proof construction and checking has a long
history \citep{deBruijn1970,BarendregtWiedijk2005}; Lean typechecker
verification supplies a contemporary example \citep{Carneiro2025Lean4Lean}.
Library maintenance addresses consistency across a growing development
\citep{vanDoornEbnerLewis2020}. Interpreting the mathematics additionally
requires reading the formal objects and hypotheses, the task pursued
by the case studies here.
